% Chương 2
\chapter{Kỹ thuật xử lý CNN}
\label{ch:chuong2}
\section{Tổng quan về mạng nơ-ron tích chập (CNN)} % TODO: bản Word thiếu tiêu đề mục 2.1, tiêu đề này do script thêm để giữ đúng số 2.1.x

\subsection{Bài toán phân loại ảnh da liễu}

Phân loại ảnh da liễu là bài toán thị giác máy tính nhằm xác định lớp bệnh tương ứng với hình ảnh tổn thương da đầu vào. Với sự phát triển của học sâu, mạng nơ-ron tích chập (Convolutional Neural Network – CNN) được sử dụng rộng rãi nhờ khả năng tự động học các đặc trưng trực quan từ dữ liệu ảnh.

Khác với phương pháp xử lý ảnh truyền thống yêu cầu thiết kế thủ công các đặc trưng, CNN học biểu diễn phân cấp thông qua quá trình huấn luyện. Các lớp đầu mạng thường học những đặc trưng cục bộ như cạnh, đường biên và biến đổi màu sắc; các lớp sâu hơn kết hợp những đặc trưng này thành biểu diễn có tính ngữ nghĩa cao hơn.

Trong đồ án này, ảnh da liễu được đưa vào mô hình AdaptiveCNN để dự đoán một trong các lớp bệnh thuộc tập dữ liệu huấn luyện. Mô hình thực hiện việc trích xuất đặc trưng, tổng hợp thông tin không gian và phân loại thông qua lớp đầu ra. Độ tin cậy của dự đoán phụ thuộc vào chất lượng dữ liệu, khả năng biểu diễn của kiến trúc và quá trình tối ưu mô hình.

\subsection{Biểu diễn ảnh dưới dạng tensor}

Ảnh màu thường được biểu diễn dưới dạng tensor ba chiều gồm số kênh màu, chiều cao và chiều rộng. Với ảnh RGB có kích thước $H\times W$, tensor ảnh có dạng:
\begin{equation}
  X\in\mathbb{R}^{3\times H\times W}
\end{equation}

Khi xử lý theo lô (batch), dữ liệu có dạng:
\begin{equation}
  X\in\mathbb{R}^{N\times C\times H\times W}
\end{equation}

Trong đó, $N$ là số ảnh trong một batch, $C$ là số kênh, $H$ là chiều cao và $W$ là chiều rộng.

Trước khi đưa vào mạng, ảnh thường được chuyển về kích thước thống nhất và chuẩn hóa giá trị pixel. Việc chuẩn hóa giúp dữ liệu có thang giá trị phù hợp với quá trình tối ưu. Các phép biến đổi tăng cường dữ liệu có thể được sử dụng trong giai đoạn huấn luyện nhằm cải thiện khả năng khái quát hóa.

\subsection{Phép tích chập}

Phép tích chập (convolution) là thành phần cốt lõi của CNN. Một kernel có kích thước xác định được áp dụng trên các vùng cục bộ của đầu vào để tạo ra feature map. Các trọng số kernel được cập nhật trong quá trình huấn luyện để học những đặc trưng hữu ích cho bài toán.

Dạng tổng quát của phép tích chập hai chiều có thể biểu diễn như sau:
\begin{equation}
  Y_{i,j}=\sum_{c=1}^{C_{\mathrm{in}}}\sum_{u=0}^{K_h-1}\sum_{v=0}^{K_w-1}W_{c,u,v}X_{c,i+u,j+v}+b
\end{equation}

Trong đó, $X$ là tensor đầu vào, $W$ là trọng số kernel, $b$ là bias và $Y$ là feature map đầu ra. Công thức được viết ở dạng đơn giản, chưa thể hiện đầy đủ các trường hợp có stride, padding hoặc dilation.

Kích thước kernel quyết định phạm vi không gian được xét trong mỗi phép tính. Kernel $3\times3$ tập trung vào vùng lân cận nhỏ, trong khi kernel $5\times5$ thu nhận thông tin từ vùng rộng hơn. Việc sử dụng nhiều kích thước kernel là một hướng tiếp cận để khai thác đặc trưng ở nhiều phạm vi không gian.

\subsection{Hàm kích hoạt và pooling}

Sau phép tích chập, hàm kích hoạt phi tuyến cho phép mạng học những quan hệ phức tạp hơn so với phép biến đổi tuyến tính đơn thuần. Một hàm phổ biến là ReLU:
\begin{equation}
  \mathrm{ReLU}(x)=\max(0,x)
\end{equation}

Pooling thực hiện giảm kích thước không gian của feature map hoặc tổng hợp thông tin từ các vùng lân cận. Max Pooling lấy giá trị lớn nhất trong mỗi vùng, trong khi Global Average Pooling (GAP) tính trung bình các giá trị trên toàn bộ không gian của từng kênh.

Các phép toán này góp phần xây dựng biểu diễn đặc trưng phù hợp cho phân loại ảnh. Trong AdaptiveCNN, việc tổng hợp feature map ở giai đoạn cuối được thực hiện bằng Adaptive Generalized Mean Pooling, được trình bày tại Mục 2.5.

\section{Mạng ResNet-18}

\subsection{Kiến trúc tổng quan}

ResNet (Residual Network) là kiến trúc CNN sử dụng kết nối tắt nhằm hỗ trợ quá trình huấn luyện mạng sâu. ResNet-18 là một biến thể gồm 18 lớp có trọng số theo cách đếm thông dụng, với các khối tích chập được tổ chức theo nhiều giai đoạn.

Kiến trúc ResNet-18 tiêu chuẩn gồm lớp tích chập ban đầu, các giai đoạn residual block, bước tổng hợp đặc trưng và lớp phân loại. Trong mỗi giai đoạn, mạng học các biểu diễn đặc trưng ngày càng trừu tượng thông qua tích chập và thay đổi kích thước không gian.

ResNet-18 được lựa chọn làm nền tảng cho AdaptiveCNN vì cung cấp cơ chế trích xuất đặc trưng phân cấp và kết nối dư, đồng thời có quy mô tương đối gọn so với nhiều kiến trúc sâu hơn.

\subsection{Residual Learning và Skip Connection}

Trong mạng CNN truyền thống, các lớp liên tiếp phải học trực tiếp ánh xạ từ đầu vào sang đầu ra. Khi mạng sâu hơn, việc tối ưu có thể gặp khó khăn do gradient suy giảm hoặc quá trình học không hiệu quả.

ResNet giải quyết vấn đề này bằng cách học phần dư (residual mapping). Đầu ra của một residual block được biểu diễn như sau:
\begin{equation}
  y=F(x,W)+x
\end{equation}

Trong đó:
\begin{itemize}
  \item $x$ là đầu vào của khối.
  \item $F(x,W)$ là phép biến đổi được học bởi các lớp bên trong khối.
  \item $W$ là tập tham số của các lớp đó.
  \item $y$ là đầu ra sau khi cộng với nhánh tắt.
\end{itemize}

Skip connection tạo đường truyền trực tiếp cho thông tin từ đầu vào đến đầu ra. Khi kích thước tensor giữa hai nhánh không tương thích, kiến trúc có thể sử dụng phép chiếu để điều chỉnh số kênh hoặc kích thước đặc trưng.

Cơ chế này giúp tối ưu mạng sâu thuận lợi hơn và là nền tảng quan trọng của AdaptiveCNN.

\subsection{Vai trò của ResNet-18 trong AdaptiveCNN}

Trong đồ án, ResNet-18 được sử dụng làm backbone để trích xuất đặc trưng từ ảnh da liễu. Tuy nhiên, AdaptiveCNN không chỉ sử dụng nguyên trạng kiến trúc ResNet-18 mà còn tích hợp cơ chế AMKC và Adaptive GeM nhằm thay đổi cách trích xuất và tổng hợp đặc trưng.

AMKC hướng đến việc kết hợp thông tin từ các kernel có kích thước khác nhau, trong khi Adaptive GeM điều chỉnh phép tổng hợp đặc trưng không gian bằng tham số $p$. Hai thành phần này được trình bày chi tiết ở các mục tiếp theo.

\section{Cơ chế tích chập đa kernel AMKC}

Các nghiên cứu về tích chập đa kernel, đặc trưng đa tỷ lệ và tích chập thích nghi cho thấy việc mở rộng cách biểu diễn đặc trưng là một hướng tiếp cận đáng chú ý trong thị giác máy tính và xử lý ảnh y khoa. Tuy nhiên, các nghiên cứu có thể khác nhau về nhiệm vụ, kiến trúc và cơ chế kết hợp đặc trưng; vì vậy cần phân biệt giữa phân loại ảnh và phân đoạn ảnh khi đối chiếu kết quả.

\subsection{Tích chập đa kernel và đa tỷ lệ}

Khan và cộng sự đề xuất MKIS-Net, một mạng nhẹ sử dụng thiết kế multi-kernel để tạo receptive field hiệu quả cho bài toán phân đoạn ảnh y khoa. Nghiên cứu đánh giá mô hình trên nhiều nhiệm vụ, bao gồm phân đoạn tổn thương da, mạch máu võng mạc và ảnh X-quang ngực. Kết quả cho thấy thiết kế đa kernel có thể hỗ trợ xây dựng mạng phân đoạn gọn nhẹ mà vẫn duy trì hiệu quả cạnh tranh [1].

Han và cộng sự đề xuất DM-CNN cho bài toán phân loại ảnh y khoa, trong đó có dữ liệu da liễu. Mô hình sử dụng mô-đun Dynamic Multi-scale Feature Fusion để lựa chọn kernel dựa trên đặc trưng ở các tầng khác nhau. Ngoài ra, kiến trúc còn kết hợp cơ chế attention có xét đến độ bất định, phương pháp pooling đa tỷ lệ và hàm mất mát đa mục tiêu [2].

So với thiết kế chỉ sử dụng một kích thước kernel cố định, các phương pháp đa kernel hoặc đa tỷ lệ hướng đến việc khai thác thông tin ở nhiều phạm vi không gian. Đây là cơ sở tham khảo cho việc nghiên cứu cách kết hợp đặc trưng từ kernel (3x3) và (5x5) trong AdaptiveCNN.

\subsection{Tích chập thích nghi và lựa chọn kernel}

Hu và cộng sự đề xuất Domain and Content Adaptive Convolution (DCAC) nhằm cải thiện khả năng tổng quát hóa giữa các miền dữ liệu y khoa. Phương pháp gồm Domain Adaptive Convolution, điều chỉnh phép tích chập dựa trên thông tin miền, và Content Adaptive Convolution, điều chỉnh theo đặc trưng toàn cục của ảnh đầu vào. Nghiên cứu tập trung vào phân đoạn ảnh y khoa giữa các miền dữ liệu khác nhau [3].

Chen và cộng sự phát triển ESKNet cho bài toán phân đoạn khối u vú trên ảnh siêu âm. Mô-đun Enhanced Selective Kernel Convolution kết hợp biểu diễn từ nhiều vùng đặc trưng và tái hiệu chỉnh thông tin theo chiều kênh lẫn không gian. Cơ chế này giúp mạng lựa chọn các đặc trưng có đóng góp cao hơn cho nhiệm vụ phân đoạn [4].

Các nghiên cứu trên cho thấy cần phân biệt ba khái niệm. Multi-kernel sử dụng nhiều kích thước kernel; multi-scale tập trung khai thác đặc trưng ở nhiều phạm vi hoặc cấp độ; còn adaptive/selective convolution có thêm cơ chế điều chỉnh hoặc lựa chọn đặc trưng dựa trên đầu vào. Những khái niệm này có liên quan nhưng không đồng nhất.

\subsection{Kết hợp convolution và attention}

AMCANet là nghiên cứu về kiến trúc convolution kết hợp attention cho bài toán khôi phục độ phân giải ảnh. Mô hình sử dụng các nhánh convolution attention với những mức lấy mẫu khác nhau nhằm kết hợp thông tin cục bộ và ngữ cảnh không gian, sau đó sử dụng cơ chế chọn lọc đặc trưng [5].

Do nhiệm vụ chính của AMCANet là super-resolution, nghiên cứu này chỉ được xem là tài liệu tham khảo gián tiếp về thiết kế convolution và attention, không phải bằng chứng trực tiếp về hiệu quả trong phân loại bệnh da liễu.

\subsection{Nhận xét và định hướng của đồ án}

Các nghiên cứu được khảo sát thể hiện những hướng tiếp cận khác nhau trong khai thác đặc trưng: thiết kế nhiều kernel, kết hợp đặc trưng đa tỷ lệ, lựa chọn kernel thích nghi và kết hợp convolution với attention. Tuy nhiên, các phương pháp này được xây dựng cho những nhiệm vụ và kiến trúc khác nhau, do đó không thể kết luận rằng một cơ chế sẽ luôn hiệu quả hơn cơ chế khác khi chuyển sang phân loại ảnh da liễu.

Từ cơ sở đó, đồ án sử dụng ResNet-18 làm backbone và tích hợp AMKC để kết hợp thông tin từ các nhánh kernel khác nhau. Adaptive GeM được sử dụng để tổng hợp feature map trước khi phân loại. Hiệu quả của thiết kế này cần được kiểm chứng bằng thực nghiệm so sánh với mô hình cơ sở và các thí nghiệm loại bỏ từng thành phần.

\section{Cơ chế tích chập đa kernel AMKC}

\subsection{Động cơ thiết kế}

Trong ảnh da liễu, các tổn thương có thể khác nhau về kích thước, hình dạng, đường biên và cấu trúc sắc tố. Vì vậy, đặc trưng thu được từ một phạm vi không gian duy nhất có thể chưa biểu diễn đầy đủ mọi dạng tổn thương.

Kernel $3\times3$ thu nhận thông tin từ vùng lân cận nhỏ, trong khi kernel $5\times5$ mở rộng vùng tiếp nhận không gian. Việc kết hợp hai nhánh cho phép mô hình đồng thời khai thác đặc trưng ở các phạm vi khác nhau mà không cần tăng độ sâu của toàn bộ mạng.

AMKC (Adaptive Multi-Kernel Convolution) trong AdaptiveCNN được thiết kế nhằm kết hợp đầu ra từ các nhánh convolution với kích thước kernel khác nhau thông qua các hệ số điều chỉnh có thể học.

\subsection{Biểu diễn toán học}

Gọi $x$ là tensor đầu vào của mô-đun. Hai nhánh tích chập được ký hiệu là:
\begin{equation}
  F_3(x)=\mathrm{Conv}_{3\times3}(x)
\end{equation}
\begin{equation}
  F_5(x)=\mathrm{Conv}_{5\times5}(x)
\end{equation}

Trong biến thể AMKC được mô tả trong đồ án, đầu ra có thể được biểu diễn dưới dạng:
\begin{equation}
  y=x+\lambda\left[\tanh(a)F_3(x)+\tanh(b)F_5(x)\right]
\end{equation}

Trong đó:
\begin{itemize}
  \item $x$ là đầu vào của mô-đun.
  \item $F_3(x)$ và $F_5(x)$ là đặc trưng từ hai nhánh kernel.
  \item $a$ và $b$ là các tham số điều chỉnh có thể học.
  \item $\tanh(\cdot)$ giới hạn hệ số điều chỉnh trong khoảng $(-1,1)$.
  \item $\lambda$ điều khiển mức đóng góp của phần đặc trưng bổ sung.
  \item $y$ là đầu ra của mô-đun.
\end{itemize}

Công thức trên thể hiện nguyên lý kết hợp đặc trưng đa kernel với một nhánh dư. Để phép cộng được thực hiện, các tensor ở hai nhánh phải có kích thước tương thích. Trong triển khai thực tế, điều này phụ thuộc vào padding, stride và cách tổ chức các lớp convolution.

Khi các nhánh convolution được khởi tạo với trọng số bằng không, phần đóng góp của chúng ban đầu bằng không. Trong điều kiện kích thước đầu vào và đầu ra tương thích, đầu ra mô-đun ban đầu gần với ánh xạ đồng nhất (y=x). Cách khởi tạo này cho phép mô hình bắt đầu từ biểu diễn gần với đường truyền ban đầu rồi học dần mức đóng góp của các nhánh bổ sung.

\subsection{Cơ chế kết hợp đặc trưng}

AMKC kết hợp thông tin từ hai phạm vi kernel bằng tổng có trọng số, thay vì chỉ nối các feature map theo chiều kênh. Các hệ số điều chỉnh cho phép mô hình học mức đóng góp của từng nhánh trong quá trình tối ưu.

So với phép cộng trực tiếp hai nhánh với trọng số cố định, các hệ số có thể học tạo thêm khả năng điều chỉnh trong huấn luyện. Tuy nhiên, đây không phải cơ chế lựa chọn kernel động theo từng ảnh nếu các hệ số chỉ là tham số học được dùng chung cho mọi đầu vào. Cần phân biệt thiết kế này với các mô hình dynamic convolution có trọng số kernel được tạo hoặc điều chỉnh theo đặc trưng của từng ảnh.

\subsection{Vai trò của AMKC trong AdaptiveCNN}

Trong AdaptiveCNN, AMKC được tích hợp nhằm mở rộng cách trích xuất đặc trưng so với lớp convolution đơn kernel tại vị trí được thiết kế. Backbone ResNet-18 tiếp tục thực hiện việc học đặc trưng phân cấp, còn AMKC bổ sung cơ chế kết hợp thông tin từ nhiều phạm vi không gian.

Việc sử dụng AMKC được kỳ vọng giúp mô hình biểu diễn đa dạng hơn các cấu trúc tổn thương da. Tuy nhiên, mức cải thiện thực tế cần được đánh giá bằng thực nghiệm, chẳng hạn so sánh với ResNet-18 cơ sở trong cùng điều kiện huấn luyện và sử dụng ablation study để xác định đóng góp của AMKC.

\section{Adaptive Generalized Mean Pooling (Adaptive GeM)}

\subsection{Tổng hợp đặc trưng không gian}

Sau các lớp convolution, mạng thu được tensor đặc trưng có dạng $C\times H\times W$. Để đưa tensor này vào lớp phân loại, cần tổng hợp thông tin không gian thành một vector có $C$ phần tử.

Global Average Pooling (GAP) thực hiện phép lấy trung bình trên toàn bộ vị trí không gian của mỗi kênh:
\begin{equation}
  f_c=\frac{1}{|\Omega|}\sum_{i\in\Omega}x_{c,i}
\end{equation}

Trong đó, $\Omega$ là tập các vị trí không gian và $x_{c,i}$ là giá trị đặc trưng tại vị trí $i$ của kênh $c$.

GAP có ưu điểm là đơn giản, ít tham số và giảm số chiều trước lớp phân loại. Tuy nhiên, phép lấy trung bình có thể làm giảm mức nổi bật của các giá trị đặc trưng lớn tại một số vị trí.

\subsection{Generalized Mean Pooling}

Generalized Mean Pooling (GeM) mở rộng phép tổng hợp bằng cách sử dụng tham số $p$:
\begin{equation}
  f_c=\left(\frac{1}{|\Omega|}\sum_{i\in\Omega}x_{c,i}^{p}\right)^{1/p}
\end{equation}

Trong đó, $p>0$ là tham số điều khiển dạng tổng hợp. Khi $p=1$, GeM tương đương GAP đối với đầu vào không âm. Khi $p$ tăng, phép tổng hợp có xu hướng nhấn mạnh các giá trị lớn hơn. Khi $p$ tiến tới vô hạn, GeM tiến gần phép lấy giá trị lớn nhất trong miền tổng hợp, với điều kiện đầu vào phù hợp.

GeM vì vậy tạo ra một cơ chế trung gian giữa lấy trung bình và lấy cực đại. Thay vì cố định hoàn toàn cách tổng hợp đặc trưng, mô hình có thể điều chỉnh tham số $p$ để phù hợp với mục tiêu huấn luyện.

\subsection{Vai trò của Adaptive GeM}

Trong AdaptiveCNN, Adaptive GeM được sử dụng để tổng hợp feature map trước lớp phân loại. Tham số $p$ được học trong quá trình tối ưu theo cách triển khai của mô hình.

So với GAP, Adaptive GeM cung cấp thêm một mức độ linh hoạt trong cách tổng hợp thông tin không gian. Đối với ảnh da liễu, điều này có thể hữu ích khi một số vùng đặc trưng có mức kích hoạt nổi bật hơn những vùng còn lại.

\section{Huấn luyện và tối ưu mô hình CNN}

\subsection{Hàm mất mát phân loại}

Với bài toán phân loại nhiều lớp, Cross-Entropy là hàm mất mát phổ biến. Gọi $z_k$ là logit của lớp $k$, xác suất dự đoán được tính bằng Softmax:
\begin{equation}
  p_k=\frac{e^{z_k}}{\sum_{j=1}^{K}e^{z_j}}
\end{equation}

Trong đó, $K$ là số lớp. Với nhãn đúng $y$, Cross-Entropy cho một mẫu được biểu diễn như sau:
\begin{equation}
  \mathcal{L}=-\log p_y
\end{equation}

Hàm mất mát đo mức độ khác biệt giữa dự đoán của mô hình và nhãn thực tế. Mục tiêu huấn luyện là tối thiểu hóa giá trị mất mát trên tập dữ liệu huấn luyện.

Nếu sử dụng các kỹ thuật như label smoothing hoặc Mixup, công thức mất mát thực tế sẽ được điều chỉnh tương ứng. Các kỹ thuật này cần được mô tả theo đúng cấu hình huấn luyện được sử dụng trong thực nghiệm.

\subsection{Lan truyền ngược và cập nhật tham số}

Quá trình huấn luyện CNN gồm hai giai đoạn chính: lan truyền xuôi và lan truyền ngược.

Trong lan truyền xuôi, ảnh đầu vào đi qua các lớp của mạng để tạo ra logits và tính hàm mất mát. Trong lan truyền ngược (backpropagation), đạo hàm của hàm mất mát theo các tham số được tính bằng quy tắc dây chuyền.

Với tham số $\theta$, quy tắc cập nhật tổng quát theo gradient descent là:
\begin{equation}
  \theta_{t+1}=\theta_t-\eta\nabla_\theta\mathcal{L}
\end{equation}

Trong đó, $\eta$ là learning rate và $\nabla_\theta\mathcal{L}$ là gradient của hàm mất mát theo tham số. Optimizer như AdamW hoặc SGD có thể sử dụng các quy tắc cập nhật và thành phần bổ sung khác so với gradient descent cơ bản.

\subsection{Learning rate và bộ tối ưu}

Learning rate kiểm soát mức độ thay đổi tham số trong mỗi bước cập nhật. Giá trị quá lớn có thể khiến quá trình tối ưu không ổn định, trong khi giá trị quá nhỏ có thể làm quá trình hội tụ chậm.

Các chiến lược điều chỉnh learning rate theo thời gian, chẳng hạn cosine scheduling, có thể giúp kiểm soát quá trình tối ưu. Khi sử dụng mô hình pretrained, có thể đặt learning rate khác nhau cho backbone và lớp phân loại để cân bằng giữa việc điều chỉnh đặc trưng đã học và học các tham số mới.

Cấu hình optimizer, learning rate, batch size, số epoch và scheduler cụ thể của AdaptiveCNN cần được trình bày ở chương mô tả thực nghiệm.

\subsection{Tăng cường dữ liệu và mất cân bằng lớp}

Data augmentation tạo ra các biến thể của ảnh huấn luyện thông qua những phép biến đổi phù hợp, chẳng hạn lật ảnh hoặc thay đổi hình học ở mức hợp lý. Mục tiêu là giảm sự phụ thuộc của mô hình vào những đặc điểm không ổn định của dữ liệu và cải thiện khả năng khái quát hóa.

Đối với dữ liệu bệnh da liễu, phép biến đổi cần bảo toàn những đặc điểm có ý nghĩa chẩn đoán. Những biến đổi làm thay đổi quá mức màu sắc hoặc cấu trúc tổn thương có thể gây sai lệch thông tin.

Dữ liệu phân loại bệnh thường có thể mất cân bằng giữa các lớp. Khi một số lớp có nhiều mẫu hơn đáng kể, mô hình có thể thiên về những lớp phổ biến. Một cách xử lý là lấy mẫu có trọng số (weighted sampling), trong đó xác suất lấy mẫu được điều chỉnh để các lớp ít mẫu có cơ hội xuất hiện nhiều hơn trong quá trình huấn luyện.

Việc sử dụng augmentation và phương pháp xử lý mất cân bằng phải được xác định theo pipeline thực tế. Những kỹ thuật không được bật trong mã nguồn không nên được mô tả như một phần của cấu hình thực nghiệm đã triển khai.

\section{Các chỉ số đánh giá mô hình}

Đánh giá mô hình phân loại ảnh da liễu cần xem xét cả khả năng dự đoán đúng tổng thể và hiệu quả trên từng lớp. Các chỉ số nên được tính trên tập kiểm tra độc lập với tập được sử dụng để tối ưu và lựa chọn checkpoint.

\subsection{Accuracy}

Accuracy đo tỷ lệ mẫu được phân loại chính xác trên tổng số mẫu:
\begin{equation}
  \mathrm{Accuracy}=\frac{\sum_{k=1}^{K}TP_k}{N}
\end{equation}

Trong đó, $TP_k$ là số mẫu thuộc lớp $k$ được dự đoán đúng và $N$ là tổng số mẫu kiểm tra.

Accuracy dễ diễn giải nhưng có thể che khuất hiệu quả thấp trên những lớp có số lượng mẫu ít.

\subsection{Precision và Recall}

Đối với lớp $k$, Precision được xác định bởi:
\begin{equation}
  \mathrm{Precision}_k=\frac{TP_k}{TP_k+FP_k}
\end{equation}

Recall được xác định bởi:
\begin{equation}
  \mathrm{Recall}_k=\frac{TP_k}{TP_k+FN_k}
\end{equation}

Trong đó, $FP_k$ là số mẫu bị dự đoán nhầm thành lớp $k$, còn $FN_k$ là số mẫu thực sự thuộc lớp $k$ nhưng bị dự đoán thành lớp khác.

Precision phản ánh mức độ chính xác của các dự đoán thuộc một lớp, trong khi Recall phản ánh khả năng phát hiện các mẫu thực sự thuộc lớp đó. Hai chỉ số này có ý nghĩa bổ sung cho nhau trong đánh giá mô hình nhiều lớp.

\subsection{F1-score và Macro-F1}

F1-score là trung bình điều hòa của Precision và Recall:
\begin{equation}
  \mathrm{F1}_k=2\,\frac{P_kR_k}{P_k+R_k}
\end{equation}

Trong đó, $P_k$ và $R_k$ lần lượt là Precision và Recall của lớp $k$.

Macro-F1 được tính bằng trung bình F1-score của tất cả các lớp:
\begin{equation}
  \mathrm{Macro\text{-}F1}=\frac{1}{K}\sum_{k=1}^{K}\mathrm{F1}_k
\end{equation}

Khác với các chỉ số trung bình có trọng số theo số mẫu, Macro-F1 coi mỗi lớp có mức đóng góp như nhau. Do đó, chỉ số này phù hợp để đánh giá khả năng phân loại khi dữ liệu có phân bố lớp không đồng đều.

\subsection{ROC-AUC}

ROC-AUC đo khả năng phân biệt giữa các lớp dựa trên sự thay đổi ngưỡng quyết định. Trong bài toán nhiều lớp, ROC-AUC có thể được tính theo chiến lược one-vs-rest, trong đó từng lớp lần lượt được xem là lớp dương và các lớp còn lại là lớp âm.

Khi báo cáo ROC-AUC, cần nêu rõ phương pháp tổng hợp đa lớp và cách tính được sử dụng. ROC-AUC bổ sung thông tin về khả năng phân biệt của điểm số dự đoán, nhưng không thay thế hoàn toàn các chỉ số đánh giá tại ngưỡng phân loại cụ thể.

\subsection{Confusion Matrix}

Confusion matrix biểu diễn số lượng mẫu theo lớp thực tế và lớp dự đoán. Các phần tử trên đường chéo thể hiện số mẫu dự đoán đúng, còn các phần tử ngoài đường chéo thể hiện những dạng nhầm lẫn giữa các lớp.

Ma trận này giúp xác định những cặp bệnh mà mô hình dễ nhầm lẫn, từ đó hỗ trợ phân tích nguyên nhân và đề xuất hướng cải tiến.

\subsection{Độ phức tạp tính toán}

Bên cạnh hiệu năng phân loại, quy mô và chi phí tính toán của mô hình cũng cần được xem xét. Số lượng tham số (Parameters) phản ánh quy mô các tham số có thể học; FLOPs hoặc MACs ước lượng lượng phép tính; thời gian suy luận phản ánh tốc độ xử lý trong điều kiện đo cụ thể.

Các chỉ số này cần được so sánh trong điều kiện nhất quán về kích thước ảnh, thiết bị và phương pháp đo. Một mô hình có độ chính xác cao hơn chưa chắc phù hợp hơn nếu chi phí tính toán tăng đáng kể.

\section{Kết luận chương}

Chương 2 đã trình bày những cơ sở lý thuyết và kỹ thuật liên quan đến việc xây dựng mô hình phân loại ảnh da liễu dựa trên CNN. Trước hết, các khái niệm về tensor ảnh, phép tích chập, hàm kích hoạt và pooling được giới thiệu làm nền tảng cho quá trình trích xuất đặc trưng.

Tiếp theo, kiến trúc ResNet-18 và cơ chế Residual Learning được trình bày nhằm làm rõ vai trò của backbone trong việc học đặc trưng phân cấp. Các nghiên cứu về tích chập đa kernel, đa tỷ lệ, tích chập thích nghi và convolution kết hợp attention cung cấp bối cảnh khoa học cho việc thiết kế các thành phần mở rộng.

Trên cơ sở đó, chương đã mô tả nguyên lý AMKC và Adaptive GeM trong AdaptiveCNN. AMKC kết hợp đặc trưng từ các nhánh kernel có kích thước khác nhau, trong khi Adaptive GeM điều chỉnh cách tổng hợp đặc trưng không gian thông qua tham số $p$. Chương cũng trình bày các nguyên tắc huấn luyện, tối ưu và các chỉ số đánh giá mô hình.

Những cơ sở này là nền tảng để trình bày kiến trúc triển khai, quy trình thực nghiệm và đánh giá hiệu quả của AdaptiveCNN trong chương tiếp theo.
